\documentclass[../../main.tex]{subfiles}
\begin{document}

{\bf [解]}
$a,b,p>0$で考える．$t=\dfrac{b}{a}$（$t>0$）とおいて$C=A-B$としてこの符号を調べることで$A,B$の大小関係を調べる．
題意より
\begin{align}
C=A-B=a^p\left[(1+t)^p-2^{p-1}(1+t^p)\right]
\end{align}
であり，$a>0$よりこの符号は
\begin{align}
 g(t)=(1+t)^p-2^{p-1}(1+t^p)
\end{align}
の符号に等しい．そこで以下$g(t)$の振る舞いを調べよう．
一階微分は
\begin{align}
g'(t)=p(1+t)^{p-1}-p\cdot2^{p-1}t^{p-1}=p\left[(1+t)^{p-1}-(2t)^{p-1}\right]
\end{align}
であるから，$p$の値に応じて以下のようになる．

\subparagraph{$1^\circ$ $0<p<1$の時}

$p-1<0$だから$x^{p-1}$は$x>0$で単調減少である．$0<t<1$では$1+t>2t$だから$(1+t)^{p-1}<(2t)^{p-1}$だから$g'(t)<0$で，$t>1$では$1+t<2t$だから$g'(t)>0$となる．
また$t=1$で$g'(1)=0$である．従って増減表は以下のようになる．ただし
\begin{align}
 g(1)=2^p-2^{p-1}\cdot2=0
\end{align}
である．

\begin{table}[htb]
  \centering
\caption{$C$の増減表}
\begin{tabular}{c|ccccc}
$t$ & $0$ & & $1$ & & \\
\hline
$g'$ & & $-$ & $0$ & $+$ & \\
\hline
$C$ & & $\searrow$ & $0$ & $\nearrow$ &
\end{tabular}
\end{table}

従って$C \ge 0$すなわち$A \ge B$である．等号が成立するのは$t=1$つまり$a=b$の時である．

\subparagraph{$2^\circ$ $p=1$の時}

この時は$A,B$の表式から明らかに$A=B$である．

\subparagraph{$3^\circ$ $1<p$の時}

$p-1>0$だから$x^{p-1}$は単調増加で，$1^\circ$と逆に$0<t<1$で$1+t>2t$より$(1+t)^{p-1}>(2t)^{p-1}$より$g'(t)>0$で，$t>1$では$g'(t)<0$となる．
よって増減表は以下のようになる．

\begin{table}[htb]
  \centering
\caption{$C$の増減表}
\begin{tabular}{c|ccccc}
$t$ & $0$ & & $1$ & & \\
\hline
$g'$ & & $+$ & $0$ & $-$ & \\
\hline
$C$ & & $\nearrow$ & $0$ & $\searrow$ &
\end{tabular}
\end{table}

従って$C \le 0$すなわち$A\le B$で，等号成立は$t=1$つまり$a=b$の時である．
\casesend

以上をまとめて，
\begin{align}
p=1\ \text{または}\ a=b\ \text{の時}\ A=B,\qquad
0<p<1\ \text{かつ}\ a\ne b\ \text{の時}\ A>B,\qquad
1<p\ \text{かつ}\ a\ne b\ \text{の時}\ A<B
\end{align}
である．

\newpage

\end{document}
